\documentclass[11pt]{amsart}
\usepackage{amsmath,amssymb,amsthm,mathtools}
\usepackage[margin=1in]{geometry}
\usepackage{hyperref}

\newtheorem{proposition}{Proposition}
\newtheorem{corollary}[proposition]{Corollary}
\theoremstyle{remark}
\newtheorem{remark}[proposition]{Remark}

\title{The Parabola--Secant Divisor Chamber in the AM--GM Circle View}
\author{Jeremy Ebert}
\address{Franklin County, Ohio, USA}
\date{2026}

\begin{document}
\maketitle

\begin{abstract}
For the AM--GM coordinates
\[
X=\frac{x-y}{2},\qquad Y=\sqrt{xy},\qquad T=\frac{x+y}{2},
\]
the positive-factor boundaries $x=1$ and $y=1$ become the mirror parabolas
\[
Y^2=1-2X,\qquad Y^2=1+2X,
\]
while the constant-product boundary $xy=n$ becomes the horizontal secant $Y=\sqrt n$. We identify the exact curvilinear region enclosed by these three curves, compute its raw Euclidean area in the flat $(X,Y)$ circle view, and recover its original factor-plane area using the canonical transported measure $2\cos\theta\,dX\,dY$.
\end{abstract}

\section{The canonical positive-factor chamber}
For $n>1$ define
\[
\Omega_n=\{(x,y):x\ge1,\ y\ge1,\ xy\le n\}.
\]
This is the continuous region naturally bounded by the positive-lattice coordinate walls and the divisor hyperbola.

Under
\[
F(x,y)=\left(\frac{x-y}{2},\sqrt{xy}\right),
\]
the boundaries transform as follows:
\[
x=1\quad\Longrightarrow\quad Y^2=1-2X,
\]
\[
y=1\quad\Longrightarrow\quad Y^2=1+2X,
\]
and
\[
xy=n\quad\Longrightarrow\quad Y=\sqrt n.
\]
The two parabolas meet at $(X,Y)=(0,1)$, and the secant meets them at
\[
X=\pm\frac{n-1}{2},\qquad Y=\sqrt n.
\]
Thus $F(\Omega_n)$ is exactly the curvilinear chamber enclosed by the row-$1$ parabola, the column-$1$ parabola, and the secant $Y=\sqrt n$.

\begin{proposition}[Raw Euclidean chamber area]\label{prop:rawchamber}
The flat Euclidean area of the parabola--secant chamber is
\[
\boxed{
P_n:=\operatorname{Area}_{XY}(F(\Omega_n))
=\frac{n^{3/2}-3\sqrt n+2}{3}.}
\]
\end{proposition}

\begin{proof}
For $1\le Y\le\sqrt n$, solve the two parabolas for $X$:
\[
X_{\rm left}(Y)=\frac{1-Y^2}{2},\qquad
X_{\rm right}(Y)=\frac{Y^2-1}{2}.
\]
Their horizontal separation is therefore
\[
X_{\rm right}-X_{\rm left}=Y^2-1.
\]
Hence
\[
P_n=\int_1^{\sqrt n}(Y^2-1)\,dY
=\left[\frac{Y^3}{3}-Y\right]_1^{\sqrt n}
=\frac{n^{3/2}-3\sqrt n+2}{3}.
\]
\end{proof}

\begin{proposition}[Factor-plane area of the same chamber]\label{prop:factorchamber}
The ordinary factor-plane area of $\Omega_n$ is
\[
\boxed{
L_{\rm chamber}(n)
= n\log n-n+1.}
\]
\end{proposition}

\begin{proof}
For $1\le x\le n$, the vertical range is
\[
1\le y\le\frac{n}{x}.
\]
Therefore
\[
\operatorname{Area}_{xy}(\Omega_n)
=\int_1^n\left(\frac nx-1\right)dx
=n\log n-n+1.
\]
\end{proof}

\section{Exact transport by the angular measure}
Let $\theta$ be ordinary polar angle in the $(X,Y)$ half-plane measured from the positive $Y$ axis. The area-distortion companion result gives
\[
dx\,dy=2\cos\theta\,dX\,dY.
\]
Consequently the two distinct area values above are exact measures of the same transported region.

\begin{corollary}[Exact weighted recovery]\label{cor:weighted}
For the parabola--secant chamber,
\[
\boxed{
 n\log n-n+1
=\iint_{F(\Omega_n)}2\cos\theta\,dX\,dY.}
\]
Thus
\[
\boxed{
 n\log n-n+1
\quad\xrightarrow{F}\quad
\frac{n^{3/2}-3\sqrt n+2}{3}}
\]
means that the same region carries the first value under transported factor-area measure and the second under raw Euclidean circle area; the displayed numbers are not equal.
\end{corollary}

\begin{remark}[Why the chamber is canonical]
The two parabolas are not auxiliary fitting curves. They are precisely the images of the positive-factor walls $x=1$ and $y=1$. The secant is precisely the image of the divisor boundary $xy=n$. Hence $F(\Omega_n)$ is the natural continuous divisor chamber in the flat circle view.
\end{remark}

\section{The $n=11$ chamber}
For $n=11$,
\[
Y=\sqrt{11},\qquad X=\pm5
\]
at the secant endpoints. The raw flat-circle chamber area is
\[
\boxed{
P_{11}=\frac{8\sqrt{11}+2}{3}\approx9.510999,}
\]
while the ordinary factor-plane area of the same region is
\[
\boxed{
11\log11-10\approx16.376848.}
\]
The discrete divisor count $D(11)=29$ is a lattice-point/staircase quantity and is not identified with either continuous area. The exact relationship supplied here is geometric transport of the continuous positive-factor chamber.

\section{Conclusion}
The row-$1$ parabola, column-$1$ parabola, and constant-product secant form a canonical divisor chamber in the flat circle view. Its raw Euclidean area is $(n^{3/2}-3\sqrt n+2)/3$, while the canonical angular weight $2\cos\theta$ recovers exactly the factor-plane area $n\log n-n+1$. This closes the continuous area branch at the natural positive-factor boundary without introducing an additional fitted normalization.

\begin{thebibliography}{9}
\bibitem{Ebert2026Table} J.~Ebert, \emph{The Cutting Plane: Every Line Through the Multiplication Table Is a Conic Section}, 2026.
\bibitem{Ebert2026Area} J.~Ebert, \emph{Area Distortion and Exact Band Geometry of the AM--GM Cone}, 2026.
\end{thebibliography}

\end{document}
